\chapter{2008年上海卷（春）}

\section{填空题}

\begin{problem}
已知集合 \(A = \{x | x < -1\) 或 \(2 \leqslant x < 3\}\)，\(B = \{x | -2 \leqslant x < 4\}\)，则 \(A \cup B = \)  \fillinblank{}.
\end{problem}

\begin{problem}
计算： \( \lim_{n\to\infty}\frac{3^{n}+1}{3^{n+1}+2^{n}}= \)  \fillinblank{}.
\end{problem}

\begin{problem}
函数 \( f(x)=\frac{\sqrt{-x^{2}+x+6}}{x-1} \) 的定义域是 \fillinblank{} \fillinblank{}.
\end{problem}

\begin{problem}
方程 \( 2\cos\left(x-\frac{\pi}{4}\right)=1 \) 在区间 \( (0,\pi) \) 内的解是 \fillinblank{} \fillinblank{}.
\end{problem}

\begin{problem}
已知数列 \( \{a_{n}\} \) 是公差不为零的等差数列， \( a_{1}=1 \) . 若 \( a_{1} \) 、 \( a_{2} \) 、 \( a_{5} \) 成等比数列，则 \( a_{n}= \)
\end{problem}

\begin{problem}
化简: \( \cos\left(\frac{\pi}{3}+\alpha\right)+\sin\left(\frac{\pi}{6}+\alpha\right)= \)  \fillinblank{}.
\end{problem}

\begin{problem}
已知 P 是双曲线 \( \frac{x^{2}}{a^{2}} - \frac{y^{2}}{9} = 1 \) 右支上的一点，双曲线的一条渐近线方程为 3x - y = 0. 设 \( F_{1} \) 、 \( F_{2} \) 分别为双曲线的左、右焦点. 若 \( |PF_{2}| = 3 \) ，则 \( |PF_{1}| = \)  \fillinblank{}.
\end{problem}

\begin{problem}
已知一个凸多面体共有9个面，所有被长均为1，其平面展开图如图1所示，则该凸多面体的体积 \(V =\) .

图1
\end{problem}

\begin{problem}
已知无穷数列 \(\{a_{n}\}\) 前 \(n\) 项和 \(S_{n} = \frac{1}{3} a_{n} - 1\)，则数列 \(\{a_{n}\}\) 的各项和为 \fillinblank{} \fillinblank{}.
\end{problem}

\begin{problem}
古代“五行”学说认为：“物质分金、木、 土、水、火五种属性，金克木，木克土，土克水，水克火，火克金. ”将五种不同属性的物质任意排成一列，设事件 \(A\) 表示“排列中属性相克的两种物质不相邻”，则事件 \(A\) 出现的概率是（结果用数值表示） \fillinblank{}.
\end{problem}

\begin{problem}
已知 \(a_1, a_2, \cdots, a_n; b_1, b_2, \cdots, b_n\) (\(n\) 是正整数)，令 \(L_1 = b_1 + b_2 + \cdots + b_n, L_2 = b_2 + b_3 + \cdots + b_n, \cdots, L_n = b_n\). 某人用图2分析得到恒等式: \(a_1b_1 + a_2b_2 + \cdots + a_nb_n = a_1L_1 + c_2L_2 + c_3L_3 + \cdots + c_kL_k + \cdots + c_nL_n\)，则 \(c_k =\) (\(2 \leqslant k \leqslant n\)).

图2
\end{problem}

\begin{problem}
已知 \(A(1,2)\)、\(b(3,4)\)，直线 \(l_1: x = 0, l_2: y = 0\) 和 \(l_3: x + 3y - 1 = 0\). 设 \(P_i\) 是 \(l_i (i = 1,2,3)\) 上与 \(A, B\) 两点距离平方和最小的点，则 \(\triangle P_1P_2P_3\) 的面积是 \fillinblank{} \fillinblank{}.
\end{problem}

\section{单选题}

\begin{problem}
已知向量 \(\overrightarrow{a} = (2, -3)\)，\(\overrightarrow{b} = (3, \lambda)\)，若 \(\overrightarrow{a} // \overrightarrow{b}\)，则 \(\lambda\) 等于
\end{problem}

\begin{answer}
C
\end{answer}（\quad）

\choices
    {\( \frac{2}{3}. \)}
    {\(-2\).}
    {\( -\frac{9}{2}. \)}
    {\( -\frac{2}{3}. \)}
\begin{problem}
已知椭圆 \(\frac{x^2}{10 - m} + \frac{y^2}{m - 2} = 1\)，长轴在 \(y\) 轴上. 若焦距为 4，则 \(m\) 等于
\end{problem}

\begin{answer}
A
\end{answer}（\quad）

\choices
    {4.}
    {5.}
    {7.}
    {8.}
\begin{problem}
已知函数 \( f(x) \)、\( g(x) \) 定义在 \( \R \) 上，\( h(x) = f(x) \cdot g(x) \)，则“\( f(x) \)、\( g(x) \) 均为奇函数”是
\end{problem}

\begin{problem}
已知 \(z \in \mathbf{C}\)，且 \(|z - 2 - 2\i| = 1\)，i 为虚数单位，则 \(|z + 2 - 2\i|\) 的最小值是
\end{problem}

\begin{answer}
B
\end{answer}（\quad）

\choices
    {2.}
    {3.}
    {4.}
    {5.}
\section{解答题}

\begin{problem}
已知 \(\cos \theta = -\frac{\sqrt{2}}{3}, \theta \in \left(\frac{\pi}{2}, \pi\right)\)，求 \(\frac{2}{\sin 2\theta} - \frac{\cos \theta}{\sin \theta}\) 的值.
\end{problem}

\begin{problem}
在平面直角坐标系 \(xOy\) 中，\(A\)、\(B\) 分别为直线 \(x + y = 2\) 与 \(x\)、\(y\) 轴的交点，\(C\) 为 \(AB\) 的中点. 若抛物线 \(y^2 = 2px (p > 0)\) 过点 \(C\)，求焦点 \(F\) 到直线 \(AB\) 的距离. 
\end{problem}

\begin{problem}
已知函数 \(f(x) = \log_2(2^x + 1)\). (1) 求证: 函数 \( f(x) \) 在 \( (-\infty, +\infty) \) 内单调递增; (2) 记 \( f^{-1}(x) \) 为函数 \( f(x) \) 的反函数. 若关于 x 的方程 \( f^{-1}(x) = m + f(x) \) 在 [1, 2] 上有解，求 m 的取值范围.
\end{problem}

\begin{problem}
某厂根据市场需求开发折叠式小凳(如图3所示). 凳面为三角形的尼龙布，凳脚为三根钢管. 考虑到钢管的受力和人的舒适度等因素，设计小凳应满足: \circled{1} 凳子高度为30cm, \circled{2} 三根细钢管相交处的节点O与凳面三角形ABC重心的连线垂直于凳面和地面.

图3 (1)若凳面是边长为 \(20\mathrm{cm}\) 的正三角形，三只 凳脚与地面所成的角均为 \(45^{\circ}\)，确定节点 \(O\) 分细钢管上下两段的比值(精确到0.01); (2) 若凳面是顶角为 \( 120^{\circ} \) 的等腰三角形，腰长为 24cm，节点 O 分细钢管上下两段之比为 2:3. 确定三根细钢管的长度(精确到 0.1cm).
\end{problem}

\begin{problem}
在直角坐标平面 \(xOy\) 上的一列点 \(A_{1}(1, a_{1})\)，\(A_{2}(2, a_{2})\)，\(\cdots\)，\(A_{n}(n, a_{n})\)，\(\cdots\)，简记为 \(\{A_{n}\}\). 若由 \(b_{n} = \overline{A_{n}A_{n+1}} \cdot \overline{j}^{2}\) 构成的数列 \(\{b_{n}\}\) 满足 \(b_{n+1} > b_{n}, n = 1, 2, 3, \cdots\)，其中 \(\overline{j}^{2}\) 为方向与 \(y\) 轴正方向相同的单位向量，则称 \(\{A_{n}\}\) 为 \(T\) 点列. (1) 判断 \( A_{1}(1,1) \) ，\( A_{2}\left(2,\frac{1}{2}\right) \) ，\( A_{3}\left(3,\frac{1}{3}\right) \) ，\( \cdots \) ，\( A_{n}\left(n,\frac{1}{n}\right) \) ，\( \cdots \) ，是否为 T 点列，并说明理由; (2) 若 \( \{A_{n}\} \) 为 T 点列，且点 \( A_{2} \) 在点 \( A_{1} \) 的右上方. 任取其中连续三点 \( A_{k} \) 、 \( A_{k+1} \) 、 \( A_{k+2} \) ，判断 \( \triangle A_{k}A_{k+1}A_{k+2} \) 的形状（锐角三角形、直角三角形、钝角三角形），并予以证明； (3) 若 \( \{A_{n}\} \) 为 T 点列，正整数 \( 1 \leqslant m < n < p < q \) 满足 \( m + q = n + p \) ，求证: \[ \overline {{{A _{n} A _{q}}}} \cdot \overline {{{j}}} > \overline {{{A _{m} A _{p}}}} \cdot \overline {{{j}}}. \]
\end{problem}

\begin{problem}
已知 \(z\) 是实系数方程 \(x^{2} + 2bx + c = 0\) 的虚根，记它在直角坐标平面上的对应点为 \(P_{x}(\mathrm{Re}z,\mathrm{Im}z)\) (1) 若 \( (b, c) \) 在直线 \( 2x + y = 0 \) 上，求证： \( P_{z} \) 在圆 \( C_{1}: (x - 1)^{2} + y^{2} = 1 \) 上； (2) 给定圆 \( C: (x - m)^{2} + y^{2} = r^{2} (m, r \in \R, r > 0) \) ，则存在唯一的线段 s 满足：\circled{1} 若 \( P_{z} \) 在圆 C 上，则 \( (b, c) \) 在线段 s 上；\circled{2} 若 \( (b, c) \) 是线段 s 上一点（非端点），则 \( P_{z} \) 在圆 C 上. 写出线段 s 的表达式，并说明理由； (3)由(2)知线段s与C圆之间确定了一种对应关系，通过这种对应关系的研究，填写表一(表中 \( s_{1} \) 是(1)中圆 \( C_{1} \) 的对应线段). 表一 线段s与线段s1的关系 m、r的取值或表达式 s所在直线平行于s1所在直线 s所在直线平分线段s1 线段s与线段s1长度相等

\end{problem}

\begin{problem}
(第1至12题)
\end{problem}

\begin{problem}
\(\{x|x < 4\}\). 2. \(\frac{1}{3}\). 3. \([-2,1) \cup (1,3]\).
\end{problem}

\begin{problem}
\(x = \frac{7\pi}{12}\). 5. \(a_{n} = 2n - 1\). 6. \(\cos \alpha\).
\end{problem}

\begin{problem}
\(8.1 + \frac{\sqrt{2}}{6}\). 9.\(-1\).
\end{problem}

\begin{problem}
\(\frac{1}{12}\). 11. \(a_{k} - a_{k - 1}\). 12. \(\frac{3}{2}\).
\end{problem}

\begin{problem}
(第13至16题) 题号 13 14 15 16 代号 C D A B
\end{problem}

\begin{problem}
(第17至22题)
\end{problem}

\begin{problem}
[解]原式 \(= \frac{2}{2\sin\theta\cos\theta} -\frac{\cos\theta}{\sin\theta}\) \[ = \frac {1 - \cos^{2} \theta}{\sin \theta \cos \theta} = \frac {\sin \theta}{\cos \theta}. \] 又 \(\cos \theta = -\frac{\sqrt{2}}{3},\theta \in \left(\frac{\pi}{2},\pi\right)\)， \[ \text{所以} \sin \theta = \sqrt {1 - \frac {2}{9}} = \frac {\sqrt {7}}{3}, \] \[ \text{所以} \frac {2}{\sin 2 \theta} - \frac {\cos \theta}{\sin \theta} = - \frac {\sqrt {1 4}}{2}. \]
\end{problem}

\begin{problem}
[解]由已知可得 \(A(2,0)\). \(B(0,2)\)，\(C(1,1)\)，解得抛物线方程为 \(y^{2} = x\). 于是焦点 \(F\left(\frac{1}{4},0\right)\). 因此 点 F 到直线 AB 的距离为 \[ \left| \frac {\frac {1}{4} + 0 - 2}{\frac {\sqrt {2}}{2}} \right| = \frac {7 \sqrt {2}}{8}. \]
\end{problem}

\begin{problem}
[证明](1)任取 \( x_{1}<x_{2} \) ，则 \[ f (x _{1}) - f (x _{2}) \] \[ = \log_{2} (2 ^{x _{1}} + 1) - \log_{2} (2 ^{x _{2}} + 1) \] \[ = \log_{2} \frac {2 ^{x _{1}} + 1}{2 ^{x _{2}} + 1}. \] \[ \text{因为} x _{1} < x _{2}, \text{所以} 0 < 2 ^{x _{1}} + 1 < 2 ^{x _{2}} + 1, \] \[ \text{所以} 0 < \frac {2 ^{x _{1}} + 1}{2 ^{x _{2}} + 1} < 1, \log_{2} \frac {2 ^{x _{1}} + 1}{2 ^{x _{2}} + 1} < 0, \] \[ \text{所以} f (x _{1}) < f (x _{2}), \] 即函数 \(f(x)\) 在 \((-\infty, +\infty)\) 内单调递增. [解] (2) \(\text{因为} f^{-1}(x) = \log_2(2^x - 1)(x > 0)\)，[解法一] \(\text{所以} m = f^{-1}(x) - f(x)\) \[ = \log_{2} (2 ^{x} - 1) - \log_{2} (2 ^{x} + 1) \] \[ = \log_{2} \frac {2 ^{x} - 1}{2 ^{x} + 1} = \log_{2} \left(1 - \frac {2}{2 ^{x} + 1}\right). \] 当 \(1 \leqslant x \leqslant 2\) 时，\(\frac{2}{5} \leqslant \frac{2}{2^x + 1} \leqslant \frac{2}{3}\)，\[ \text{所以} \frac {1}{3} \leqslant 1 - \frac {2}{2 ^{x} + 1} \leqslant \frac {3}{5}, \] 因此 m 的取值范围是 \( \left[\log_{2}\left(\frac{1}{3}\right),\log_{2}\left(\frac{3}{5}\right)\right] \) . 解法二 解方程 \( \log_{2}(2^{x}-1)=m+\log_{2}(2^{x}+1) \) ，得 \( x=\log_{2}\left(\frac{2^{m}+1}{1-2^{m}}\right) \) . \[ \text{因为} 1 \leqslant x \leqslant 2, \text{所以} 1 \leqslant \log_{2} \left(\frac {2 ^{m} + 1}{1 - 2 ^{m}}\right) \leqslant 2, \] 解得 \(\log_2\left(\frac{1}{3}\right) \leqslant m \leqslant \log_2\left(\frac{3}{5}\right)\). 因此 m 的取值范围是 \( \left[\log_{2}\left(\frac{1}{3}\right),\log_{2}\left(\frac{3}{5}\right)\right] \) .
\end{problem}

\begin{problem}
[解](1)设\(\triangle ABC\)的重心为\(H\)，连结\(OH\).由题意可得，\(BH=\frac{20\sqrt{3}}{3}\).设细钢管上下两段之比为\(\lambda\).已知凳子高度为30.则\(OH=\frac{30\lambda}{1+\lambda}\).因为节点\(O\)与凳面三角形\(ABC\)重心的连线与地面垂直，且凳面与地面平行，因此 \( \angle OBH \) 就是 OB 与平面 ABC 所成的角，亦即 \( \angle OBH = 45^{\circ} \) . \[ \text{因为} B H = O H, \text{所以} \frac {3 0 \lambda}{1 + \lambda} = \frac {2 0 \sqrt {3}}{3}, \] 解得，\( \lambda=\frac{2\sqrt{3}}{9-2\sqrt{3}}\approx0.63 \) ，即节点O分细钢管上下两段的比值约为0.63.

图4 (2) 设 \(\angle B = 120^{\circ}\)，\(\text{所以} AB = BC = 24\)，\(AC = 24\sqrt{3}\). 设 \(\triangle ABC\) 的重心为 \(H\)，则 \(BH = 8\)，\(AH = 8\sqrt{7}\)，由节点 \(O\) 分细钢管上下两段之比为 2:3，可知 \(OH = 12\). 设过点 \(A, B, C\) 的细钢管分别为 \(AA', BB'\)，\(CC'\)，则 \(AA' = CC' = \frac{5}{2}OA = \frac{5}{2}\sqrt{OH^2 + AH^2} = 10\sqrt{37} \approx 60.8\)，\[ B B ^{\prime} = \frac {5}{2} O B = \frac {5}{2} \sqrt {O H ^{2} + B H ^{2}} = 1 0 \sqrt {1 3} \approx 3 6. 1, \] 因此 对应于 A、B、C 三点的三根细钢管长度分别为 60.8cm, 361cm 和 60.8cm.
\end{problem}

\begin{problem}
解(1) \( a_{m}=\frac{1}{n},\text{所以} b_{n}=\frac{1}{n+1}-\frac{1}{n}=\frac{-1}{n(n+1)} \) ，显然 \( b_{n+1}>b_{n} \)
\end{problem}

\begin{problem}
[证明(1)由题意可得 \(2b + c = 0\) ，解方程\(x^{2} + 2bx - 2b = 0\) ，得 \(z = -b\pm \sqrt{-2b - b^2}\i\)， 因此 点 \( P_{2}(-b, \sqrt{-2b-b^{2}}) \) 或 \[ P _{z} (- b, - \sqrt {- 2 b - b ^{2}}), \] 将点 \(P_{2}\) 代入圆 \(C_1\) 的方程，等号成立， 因此 \( P_{2} \) 在圆 \( C_{1}:(x-1)^{2}+y^{2}=1 \) 上. (2)[解法一]当 \( \Delta<0 \) ，即 \( b^{2}<c \) 时，解得 \[ z = - b \pm \sqrt {c - b ^{2}} \i, \] 因此 点 \( P_{z}(-b, \sqrt{c-b^{2}}) \) 或 \( P_{z}(-b, -\sqrt{c-b^{2}}) \) . 由题意可得 \( (-b-m)^{2}+c-b^{2}=r^{2} \) ，整理后得 \( c=-2mb+r^{2}-m^{2} \) . \[ \text{因为} \Delta = 4 (b ^{2} - c) < 0, \] \[ (b + m) ^{2} + c - b ^{2} = r ^{2}, \] 因此 \( b \in (-m - r, -m + r) \) . 因此 线段 s 为 c = -2mb + \( r^{2} - m^{2} \) ，\( b \in [-m - r, -m + r] \) . 若 \( (b,c) \) 是线段s上一点(非端点)，则实系数方程为 \( x^{2}+2bx-2mb+r^{2}-m^{2}=0,\quad b\in(-m-r,-m+r) \) . 此时 \(\Delta < 0\) ，且点 \(P_{2}(-b,\sqrt{r^{2} - (b + m)^{2}})\) \(P_{2}(-b, - \sqrt{r^{2} - (b + m)^{2}})\) 在圆 \(C\) 上. [解法二]设 \(z = x + yi\) 是原方程的虚根，则\((x + yi)^{2} + 2b(x + yi) + c = 0\)， 解得 \( \left\{\begin{aligned}x&=-b\\ y^{2}&=x^{2}+2bx+c\end{aligned}\right. \) \circled{1} 由题意可得， \((x - m)^2 +y^2 = r^2\) \circled{3} 解\circled{1}、\circled{2}、\circled{3}得 \(c = -2mb + r^2 - m^2\) 以下同解法一. [解](3)表一 线段s与线段s1的关系 m、r的取值或表达式 s所在直线平行于s1所在直线 m=1,r≠1 s所在直线平分线段s1 \( r^{2}-(m-1)^{2}=1,m\neq1 \) 线段s与线段s1长度相等 \( (1+4m^{2})r^{2}=5 \)
\end{problem}

